\chapter{The Workholding menu}\label{ch:workholding_menu}

The \menu{Workholding} menu is for the chuck body and the jaws. Chapter~\ref{ch:workholding} explains them.

\begin{center}
\begin{tabularx}{\textwidth}{@{}l Y@{}}
\toprule
\textbf{Entry} & \textbf{What it does} \\
\midrule
\menu{Show workholding} & Shows the list of workholding items in the middle of the window. \\
\menu{Add workholding} & Adds a \emph{chuck body} or \emph{chuck jaws}. \\
\menu{Library} & Chooses and manages the workholding library. \\
\bottomrule
\end{tabularx}
\end{center}

\section{Show workholding}
\menupath{Workholding}{Show workholding} shows the Workholding view in the middle: the items of the library, a picture of the chosen item and the buttons \button{Add workholding}, \button{Delete} and \button{Sketch editor...}. The form of the chosen item is on the right.

The text under the picture reminds you that the jaws and the chuck body the program uses are chosen in the Setup. If none is chosen there, nothing is drawn and the checks use the flat clamping plane of the Setup.

\section{Add workholding}
\menupath{Workholding}{Add workholding} offers \emph{Chuck body} and \emph{Chuck jaws}. The new item has a placeholder shape; you reshape it in the sketch editor.

\section{Library}
The \menu{Library} submenu works as described for tools (Section~\ref{sec:tools_menu:library}). Workholding is kept in a library of its own, so you can have one library for each of your lathes or chucks.
